% Complete English Section 5 fragment; no preamble is included.
\IfFileExists{section5/experiment_values.tex}{%
  % Generated after a successful output-checked Table 11 run.
\newcommand{\ExperimentRunId}{20260810T052530Z}
\newcommand{\ExperimentHostOS}{Microsoft Windows 11 Home Chinese Edition 64-bit, version 10.0.26200}
\newcommand{\ExperimentRunnerOS}{Ubuntu 24.04.3 LTS}
\newcommand{\ExperimentKernel}{6.6.87.2-microsoft-standard-WSL2}
\newcommand{\ExperimentCPU}{Intel(R) Core(TM) i7-14700}
\newcommand{\ExperimentPhysicalCores}{20}
\newcommand{\ExperimentLogicalCPUs}{28}
\newcommand{\ExperimentPhysicalMemoryGiB}{24}
\newcommand{\ExperimentRepetitions}{30}
\newcommand{\ExperimentPython}{Python 3.12.3}
\newcommand{\ExperimentRust}{rustc 1.96.0 (ac68faa20 2026-05-25)}
\newcommand{\ExperimentCargo}{cargo 1.96.0 (30a34c682 2026-05-25)}
\newcommand{\ExperimentTFHErs}{0.8.7}
\newcommand{\TableElevenHamming}{92.46}
\newcommand{\TableElevenSymbolInterval}{277.38}
\newcommand{\TableElevenSymbolThreshold}{275.99}
\newcommand{\TableElevenMedian}{274.93}
%
}{%
  % Generated after a successful output-checked Table 11 run.
\newcommand{\ExperimentRunId}{20260810T052530Z}
\newcommand{\ExperimentHostOS}{Microsoft Windows 11 Home Chinese Edition 64-bit, version 10.0.26200}
\newcommand{\ExperimentRunnerOS}{Ubuntu 24.04.3 LTS}
\newcommand{\ExperimentKernel}{6.6.87.2-microsoft-standard-WSL2}
\newcommand{\ExperimentCPU}{Intel(R) Core(TM) i7-14700}
\newcommand{\ExperimentPhysicalCores}{20}
\newcommand{\ExperimentLogicalCPUs}{28}
\newcommand{\ExperimentPhysicalMemoryGiB}{24}
\newcommand{\ExperimentRepetitions}{30}
\newcommand{\ExperimentPython}{Python 3.12.3}
\newcommand{\ExperimentRust}{rustc 1.96.0 (ac68faa20 2026-05-25)}
\newcommand{\ExperimentCargo}{cargo 1.96.0 (30a34c682 2026-05-25)}
\newcommand{\ExperimentTFHErs}{0.8.7}
\newcommand{\TableElevenHamming}{92.46}
\newcommand{\TableElevenSymbolInterval}{277.38}
\newcommand{\TableElevenSymbolThreshold}{275.99}
\newcommand{\TableElevenMedian}{274.93}
%
}

\section{Experiments}

The algorithms and reproduction materials are prepared for open-source release
at
\url{https://github.com/3458463857-tech/A-New-Framework-for-Efficient-Multivariate-Functional-Bootstrapping.git}.
The current release includes the Section~4 compression code, the verified
Table~11 benchmark and raw records, and the auxiliary applicable-function
catalogue. In particular, the English \texttt{applicable\_functions.pdf}, its
\TeX{} source, the intermediate-route comparison program, and the corresponding
CSV/JSON outputs are released at the same address. The implementation for
Section~3 is not included in this release and will be open-sourced in a later
update.

\paragraph{Experimental environment.}
All experiments were performed on the same desktop machine, equipped with an
Intel(R) Core(TM) i7-14700 processor with 20 cores and 28 threads, a reported
base frequency of $2.10$~GHz, and 24~GB RAM. The host operating system was
Microsoft Windows~11 Home Chinese Edition 64-bit, version 10.0.26200. The final
Table~11 rerun (record \texttt{\ExperimentRunId}) was executed in
\ExperimentRunnerOS{} under WSL2, with kernel \ExperimentKernel. Its reference
programs used \ExperimentPython, while the online benchmark used
\ExperimentRust, \ExperimentCargo, and TFHE-rs~\ExperimentTFHErs. The division
experiments reported in Table~\ref{tab:div_time} were obtained earlier on the
same Windows~11 desktop. Unless otherwise stated, reported times are arithmetic
means of online homomorphic computation.

For the final Table~11 run, we performed five untimed warm-up evaluations and
then \ExperimentRepetitions{} timed evaluations on independently sampled valid
inputs. The timed region contains only online homomorphic evaluation: key
generation, LUT generation, input encryption, and output decryption are
excluded. Every result ciphertext is decrypted after the timer stops and
compared with the cleartext function value. The complete environment record,
dependency tree, representation-generation trace, and per-repetition outputs
are included in
\texttt{experiment\_records/run\_\ExperimentRunId}.

\subsection{Bootstrapping Implementation of Non-negative Division with Remainder}

Consider non-negative integer division with remainder
\[
g(m,d)=\lfloor m/d\rfloor,
\]
with domain $[t]\times([t]\setminus\{0\})$. Our method decomposes this
bivariate LUT into two univariate logarithmic functional bootstrappings, one
homomorphic subtraction, and one outer exponential functional bootstrapping.
According to the complexity analysis in Section~3.1.1, under theoretically
optimal parameters, the total equivalent cost is
\[
\left(
2+2\left\lceil\frac{\lfloor A\rfloor+2}{t}\right\rceil
\right)\times\Z[\zeta_q]\text{-BR},
\qquad
A=\frac{2\ln(t-1)}{\ln(1+\frac{1}{t-1})}.
\]

Table~\ref{tab:div_record} compares the number of lookup-table records among
our method,~\cite{BBR26}, and~\cite{OKC21}. For~\cite{BBR26}, we use the same
estimation method as in Section~4: choose $\lambda+2a$ as small as possible
subject to
\[
1.05\cdot t^\ell<(a+1)(\ell+\lambda),
\]
and let each PBS lookup table have size equal to the smallest prime $p$ greater
than $t$. Thus the record count is $(\lambda+2a)p$. For~\cite{OKC21}, their
Section~4.4 states that $N\tau$ torus values must be stored for the
division-by-constant test vectors; using their experimental parameter $N=1024$
and $\tau=t+1$, this gives $1024(t+1)$ records. For our method, the theoretically
optimal scale gives $t'=\lfloor A\rfloor+1$; two size-$t$ inner tables and the
size-$P=2t'+1$ outer plaintext interval therefore contain
\[
2t+P=2t+2\lfloor A\rfloor+3
\]
records. Our method reduces the record count in all three parameter settings.

\begin{table}[!tbp]
  \centering
  \small
  \caption{Lookup-table compression for division with remainder.}
  \label{tab:div_record}
  \begin{tabular}{c|c|c|c|c}
    \hline
    $\ell$ & $t$
    & \makecell{LUT size of\\\cite{BBR26}}
    & \makecell{LUT size of\\\cite{OKC21}}
    & \makecell{LUT size of\\our method} \\
    \hline
    $2$ & $16$ & $731$   & $17408$ & $201$ \\
    $2$ & $32$ & $3293$  & $33792$ & $499$ \\
    $2$ & $64$ & $12194$ & $66560$ & $1183$ \\
    \hline
  \end{tabular}
\end{table}

Table~\ref{tab:div_time} reports the measured online times of our implementation
and the model-based figures for~\cite{BBR26} and~\cite{OKC21}.

\begin{table}[!tbp]
  \centering
  \small
  \caption{Online computation times and model projections for division with
  remainder.}
  \label{tab:div_time}
  \begin{tabular}{c|c|c|c|c}
    \hline
    $\ell$ & $t$
    & \makecell{Model projection for\\\cite{BBR26} (ms)}
    & \makecell{Linear estimate for\\\cite{OKC21} (ms)}
    & \makecell{Our measured\\online time (ms)} \\
    \hline
    $2$ & $16$ & $79.2$   & $930.00$  & $267.02$ \\
    $2$ & $32$ & $356.8$  & $1860.00$ & $305.17$ \\
    $2$ & $64$ & $1321.3$ & $3720.00$ & $381.46$ \\
    \hline
  \end{tabular}
\end{table}

For~\cite{OKC21}, the original paper reports $0.93$~s for 4-bit division,
corresponding to $t=16$, and its Algorithm~10 enumerates all possible divisors,
so the cost scales linearly with $t$. Hence we use $930$~ms for $t=16$ and
linearly estimate $1860$~ms and $3720$~ms for $t=32$ and $t=64$, respectively.

For~\cite{BBR26}, the projection is reproducible from its Appendix Table~4 and
our optimistic dimension-bound PBS counts. These target counts minimize
$\lambda+2a$ under the stated inequality; they are not call counts of a
decomposition produced by the randomized search. At $t=16$, their 8-to-8-bit
experiment with chunk size $s=16$ takes $105$~ms for $57$ PBS calls, giving the
calibration $c_{17}=105/57=1.842$~ms per call for embedding modulus $p=17$. A
one-output division LUT requires $43$ calls, hence $43c_{17}=79.2$~ms. Because
no timings for $p=37$ or $p=67$ are reported, we use the explicit linear scaling
rule
\[
c_p=c_{17}\frac{p}{17},
\]
and multiply it by the computed call counts $89$ and $182$. This gives
$356.8$~ms and $1321.3$~ms. The rule captures the stated dependence on LUT
modulus and bootstrapping count, but not architecture, cache,
parameter-selection, or parallelization effects.

The prior-work columns in Table~\ref{tab:div_time} are therefore cross-platform
model calculations rather than controlled same-machine benchmarks. They are
included to make the estimation method explicit; a same-platform execution of
all implementations would be required for a definitive end-to-end speedup
comparison.

\subsection{Bootstrapping Results for General Multivariate Functions}

This subsection evaluates the general multivariate-function decomposition from
Section~4. Before cryptographic timing, the released Section~4 program
regenerates the four representations and exhaustively checks that equal
intermediate sums never correspond to different function values.

The Hamming-weight interval uses the identity representation $p_i=(0,1)$ and
therefore needs only the outer PBS after homomorphic addition. The two
symbol-set functions use, respectively, $p_i=(0,1,0,1)$ and
$p_i=(1,0,1,0)$. In each six-variable case, three inputs are packed as
$z=x_i+t x_{i+1}+t^2x_{i+2}$, and one inner LUT maps $z$ directly to
$p_i(x_i)+p_{i+1}(x_{i+1})+p_{i+2}(x_{i+2})$. Thus two packed inner PBS calls
are followed by one outer PBS. For the lower-median function, the deterministic
simulated-annealing search uses seed $20260810$ and returns the symmetric rows
$p_i=(7,0,5)$ with span $42$; the same three-variable packing again gives two
inner PBS calls and one outer PBS. The four representations are exhaustively
checked over $1024$, $4096$, $4096$, and $729$ inputs, respectively.

Table~\ref{tab:multi_time_record} reports the final output-checked online times
using the TFHE-rs parameter set
\texttt{PARAM\_\allowbreak MESSAGE\_6\_\allowbreak CARRY\_0\_\allowbreak
KS\_\allowbreak PBS}, whose plaintext space is
$[64]$. It accommodates the ten-input Hamming sum, each packed value
($t^3\leq64$), and every representation sum. Except for the identity case, the
timed computation contains
\[
\operatorname{PBS}(P_1),\operatorname{PBS}(P_2),
\qquad
P_1(z_1)+P_2(z_2)\ \text{by homomorphic addition},
\qquad
\operatorname{PBS}(q).
\]

For~\cite{BBR26}, the projected running times use the same transparent,
optimistic calibration principle. We divide the nearest matching total time in
its Appendix Table~4 by the corresponding $\gamma=1.05$ call count in its
Table~2, then multiply by the lower-bound-style one-output call count obtained
by minimizing under the dimension inequality. For $(t,\ell)=(2,10)$ this is
$(477/238)\cdot81=162.3$~ms; for $(4,6)$ it is
$(847/392)\cdot178=384.6$~ms; and for $(3,6)$ we use the nearest $p=5$, 10-bit
calibration, $(311/174)\cdot71=126.9$~ms. These are cross-platform model
projections, not measured baselines.

\begin{table}[!tbp]
  \centering
  \small
  \caption{Output-checked online computation times for general multivariate
  functions.}
  \label{tab:multi_time_record}
  \begin{tabular}{c|c|c|c}
    \hline
    \makecell{Function}
    & \makecell{Parameters}
    & \makecell{Model projection for\\\cite{BBR26} (ms)}
    & \makecell{Our measured\\online time (ms)} \\
    \hline
    \makecell{Hamming-weight\\interval function}
    & \makecell{$t=2$, $\ell=10$\\$[a,b]=[4,6]$}
    & $162.3$
    & $\TableElevenHamming$ \\
    \hline
    \makecell{General symbol-set\\interval function}
    & \makecell{$t=4$, $\ell=6$\\$A=\{1,3\}$, $[a,b]=[2,4]$}
    & $384.6$
    & $\TableElevenSymbolInterval$ \\
    \hline
    \makecell{General symbol-set\\threshold function}
    & \makecell{$t=4$, $\ell=6$\\$A=\{0,2\}$, $\tau=3$}
    & $384.6$
    & $\TableElevenSymbolThreshold$ \\
    \hline
    \makecell{Lower-median\\function}
    & \makecell{$t=3$, $\ell=6$}
    & $126.9$
    & $\TableElevenMedian$ \\
    \hline
  \end{tabular}
\end{table}

The measured and projected columns were produced on different platforms and
should not be read as a controlled speed comparison. The released Table~11
records establish that the four decompositions execute correctly with PBS-call
counts $1,3,3,3$ and provide every output-checked timed repetition used for the
reported means.
